\documentclass[10pt,a4paper]{amsart}
\usepackage[margin=19mm]{geometry}
\usepackage{amsmath,amssymb,mathtools}
\usepackage[scr=rsfs,cal=euler]{mathalfa}
\usepackage{xfrac}
\usepackage[unicode,hidelinks,bookmarks=false]{hyperref}
\newcommand{\ie}{i.e.\ }
\newcommand{\cf}{cf.\ }
\newcommand{\resp}{respectively\ }
\newcommand{\Subsection}{\subsection}
\providecommand{\nobreakdash}{\nobreak\hbox{-}}
\setlength{\emergencystretch}{2em}
\multlinegap0pt
\usepackage{bm}
\usepackage{bbm}
\usepackage{dsfont}
\usepackage{xspace}
\usepackage{cancel}
\usepackage{mathtools}
\usepackage{amsthm}
\makeatletter
\@ifundefined{theorem}{\newtheorem{theorem}{Theorem}}{}
\@ifundefined{proposition}{\newtheorem{proposition}[theorem]{Proposition}}{}
\makeatother
\title[Kinematic Varieties for Massless Particles]{Kinematic Varieties for Massless Particles}
\author{Smita Rajan, Svala Sverrisd\'ottir, Bernd Sturmfels}
\date{Source: arXiv:2408.16711v2 (CC BY 4.0)}
\begin{document}
\maketitle

\noindent\emph{Source and licence.} Kinematic Varieties for Massless Particles. Version arXiv:2408.16711v2 (CC BY 4.0), distributed under the Creative Commons Attribution 4.0 International licence, as recorded in the arXiv licence metadata for this exact version and shown on the abstract page https://arxiv.org/abs/2408.16711v2. Full rights notice: licenses/arxiv-2408.16711-CC-BY-4.0.txt.\par\smallskip

\noindent\textit{Editorial scope.} Counted unit: the Proposition whose source label is \texttt{prop:Cproperties} in the article (source0.tex lines 627--635, source label \texttt{prop:Cproperties}), delivered together with the article's own complete original proof of it (source0.tex lines 637--682, one \texttt{proof} environment of 46 lines). No mathematics is written, repaired or re-derived: statement and proof are the article's own text line for line. Required context retained uncounted: eq:Omega4 (lines 462--465); eq:Cproperties (lines 576--579). Every definition, display and label of those lines is retained unchanged. Omitted from this package and not redistributed: 1--461, 466--575, 580--626, 636--636, 683--1585. Nothing inside a retained excerpt is omitted or altered: every retained body is delivered exactly as the article's lines are written, so no omission receipt of any kind accompanies this package. Citations: the retained excerpts carry no citation command at all, so no bibliography entry of the article is transcribed and no reference is redistributed. Reference adaptation: none was needed and none was made; every reference inside the delivered unit resolves inside it, so no unresolved reference or citation is produced. Printed slips kept literal: no printed word, symbol, hypothesis or number of the counted statement or of its proof was altered. Layout and class: the article's own class is \texttt{ourlematema} with packages including amsthm, bm, mathtools, amsmath, amssymb, hyperref, color, MnSymbol, graphicx, caption, xypic, verbatim, chemfig, chemnum; the wrapper is the lane's pinned \texttt{amsart} editorial wrapper at 10pt on A4, which restates the theorem-like environments the retained text needs and adds the article's own macro definitions that the retained text needs (none), with extra wrapper packages (bm, bbm, dsfont, xspace, cancel, mathtools). The delivered numbering is the unit's own. Assets: no retained span contains a figure environment, a graphics-inclusion command, a PDF-page inclusion, a table or a TikZ picture; no image file is copied into this package, the package declares no asset dependency and no image file is redistributed with it. Licence: version arXiv:2408.16711v2 of this article is distributed under the Creative Commons Attribution 4.0 International licence, as recorded in the arXiv licence metadata for this exact version and shown on the abstract page https://arxiv.org/abs/2408.16711v2. A full rights notice is delivered at licenses/arxiv-2408.16711-CC-BY-4.0.txt. Characters: every retained excerpt is pure ASCII, so the delivered engine, XeLaTeX with its default Latin Modern text font, reports no missing character.
\begin{equation}
\label{eq:Omega4}
    \Gamma_{d}\,\,=\,\,-i^{k-1} \cdot \Gamma_1 \Gamma_2 \,\cdots\, \Gamma_{d - 1}.
\end{equation}

\begin{equation}
\label{eq:Cproperties}
CP = -P^TC \,\,\, \text{ if $d=2k$ is even}, \quad CP = (-1)^kP^TC  \,\,\, \text{ if $d=2k+1$ is odd}.
\end{equation}

\begin{proposition} \label{prop:Cproperties}
    The following properties hold for the charge conjugation matrix $C$ associated with particles
    in spacetimes of dimension $d=2k$ and $d=2k+1$:
    \begin{enumerate}
        \item  $C$ is symmetric when $k \equiv 0,3 \!\! \mod{4}$; otherwise it is skew symmetric.
        \item $C$ is block diagonal when $k \equiv 0 \!\!\! \mod{2}$; otherwise it is anti-block diagonal.
        \item The $2^{k - 1} \times 2^{k - 1}$ blocks of $C$ are skew symmetric when $k \equiv 2,3 \!\! \mod{4}$; otherwise the blocks are symmetric.
    \end{enumerate}  
\end{proposition}

\begin{proof}
    The charge conjugation matrix can be written as follows:
\begin{equation}
\label{eq:CCM}
\begin{matrix}
    C & = & \Gamma_{d + 1} \Gamma_4 \Gamma_6 \,\cdots\,
    \Gamma_{d-2} \Gamma_{d} \Gamma_1 \,\,\,\,\, \text{ if $d \equiv 0\! \!\!\!\mod{4}$,}
    \\
    C & = & \!\!\!\!\!\! \Gamma_4  \Gamma_6 \,\cdots\, \Gamma_{2k-2} \Gamma_{2k} \Gamma_1 \,
    \quad \text{ otherwise.}
    \end{matrix}
\end{equation}
    Here, we set $\Gamma_{d + 1} = -i^{k - 1} \Gamma_1 \cdots \Gamma_d$,
    which is analogous to (\ref{eq:Omega4}).    
    One can check that this $C$ defines an equivariant map to the dual representation and 
    (\ref{eq:Cproperties}) holds.
    
The Dirac matrices are either symmetric or skew symmetric. Namely, we~have
\begin{equation}
\label{eq:eithersymmetric}
    \Gamma_1^T\,=\, -\Gamma_1, \quad \Gamma_2^T\, =\,\Gamma_2,  \quad {\rm and} \quad
     \Gamma_{2i - 1}^T \,= \, \Gamma_{2i - 1},\,\,\,\,
     \Gamma_{2i}^T = -\Gamma_{2i} \quad  {\rm for} \,\,\, i \,\geq \,2 .
\end{equation}
First suppose 
  $d \not\equiv 0 \!\! \mod{4}$. Then (\ref{eq:eithersymmetric}) and the Clifford algebra relations imply
        $$
    C^T \,\,=\,\, (-1)^k \Gamma_1 \Gamma_{2k} \Gamma_{2k-2}\, \cdots\,\Gamma_6 \Gamma_4 
    \,\,=\,\, (-1)^k(-1)^{\frac{k(k - 1)}{2}} C \,\,= (-1)^{\frac{(k + 1)k}{2}} C.
    $$
    We conclude that the matrix  $C$ is symmetric if $k \equiv 3 \!\! \mod{4}$
    and it is skew symmetric if $k \equiv 1,2 \!\! \mod{4}$.
          Now suppose that $d \equiv 0 \!\! \mod{4}$. In this case, we find
    $$
    C^T\,\, =\,\, (-1)^k \Gamma_1 \Gamma_{d} \Gamma_{d-2}\, \cdots\, \Gamma_6 \Gamma_4 \Gamma_{d + 1} \,\,
    = \,\,(-1)^k (-1)^{\frac{(k + 1)k}{2}} C \,\,=\,\, (-1)^{\frac{(k + 3)k}{2}} C.
    $$
    Hence $C$ is symmetric if $k \equiv 0 \!\! \!\mod{4}$; otherwise it is skew symmetric. 
    Bearing in mind that $d \in \{2k,2k+1\}$, this completes
    the proof of part 1 in Proposition~\ref{prop:Cproperties}.
    
    All  Dirac matrices $\Gamma_i$ are anti-block diagonal, except the last one when $d$ is odd. Since $C$ always has $k$ anti-block diagonal terms, we find that $C$ is block diagonal if $k \equiv 0 \!\! \mod{2}$. Otherwise it is 
    anti-block diagonal. This proves part~2.

The verification of  part 3 is similar, but it is a bit more technical.
\end{proof}

\end{document}
